\begin{ouverture}[lunette]{17}{Réflexion et réfraction}{Ondes et signaux}{Observer · Modéliser · Expérimenter · Expliquer}
\objectif{Mesurer les angles par rapport à la normale.}
\objectif{Appliquer les lois de Snell-Descartes.}
\objectif{Déterminer un indice et expliquer la dispersion.}
\prerequis{Proportionnalité, conversions et lecture de graphiques.}
\end{ouverture}
\essentiel
\essSection{1}{À l’interface entre deux milieux}
Une partie de la lumière peut être réfléchie, une autre transmise avec un changement de direction : la réfraction. La normale est la droite perpendiculaire à l'interface au point d'incidence. Tous les angles des lois suivantes sont mesurés avec cette normale.
\schema[corrige]{refraction}
\begin{loi}[Réflexion]Le rayon réfléchi reste dans le plan d'incidence et l'angle de réflexion est égal à l'angle d'incidence : $i'=i$.\end{loi}\schema[corrige]{reflexion}
\essSection{2}{L’indice et la réfraction}
\begin{loi}[Snell-Descartes]Les rayons incident, réfracté et la normale sont dans le même plan ; $n_1\sin i=n_2\sin r$. L'indice $n$ est sans unité. On prend généralement $n_{\mathrm{air}}\simeq1,00$.\end{loi}
Si la lumière entre dans un milieu d'indice plus élevé, elle se rapproche de la normale : $r<i$. Pour trouver un angle, utiliser la fonction inverse du sinus et vérifier le mode degré de la calculatrice.
\begin{exemple}De l'air vers un verre d'indice 1,50 avec $i=30,0^\circ$, $\sin r=\sin30^\circ/1,50=1/3$ ; $r\simeq19,5^\circ$.\end{exemple}
\essSection{3}{Mesures et dispersion}
Pour une série de mesures air-milieu, tracer $\sin i$ en fonction de $\sin r$. Si le modèle convient, les points se regroupent près d'une droite passant par l'origine, de pente $n_{\mathrm{milieu}}/n_{\mathrm{air}}$.
\schema[corrige]{snell-graphe}
L'indice d'un matériau dispersif varie avec la longueur d'onde. Un prisme dévie alors différemment les radiations et décompose une lumière blanche. Un réseau peut également produire un spectre, par un mécanisme différent qui n'est pas détaillé ici. Un prisme ne crée pas les radiations : il les sépare.
\schema[corrige]{prisme}
\essSection{4}{Expérimenter et modéliser}
\begin{experience}[Tester une loi]Diriger un faisceau adapté vers une surface plane transparente, repérer normale et point d'incidence, mesurer $i$, $i'$ et $r$ pour plusieurs positions. Tracer les sinus, rechercher la proportionnalité et estimer l'indice. Éviter toute exposition des yeux au faisceau ou à sa réflexion.\end{experience}
\exercices
\rubrique{Prendre confiance}
\begin{blocLivre}{ex-17-1}\begin{exercice}[1]{La bonne référence angulaire}
Un rayon incident fait 35° avec la normale.
\begin{questions}
\item Donner l’angle d’incidence.
\item Donner l’angle qu’il fait avec la surface.
\end{questions}
\end{exercice}\end{blocLivre}
\begin{blocLivre}{ex-17-2}\begin{exercice}[1]{Réfléchir un rayon}
L’angle d’incidence vaut 40°.
\schema{reflexion}
\begin{questions}
\item Donner l’angle de réflexion.
\item Situer le rayon réfléchi par rapport à la normale.
\end{questions}
\end{exercice}\end{blocLivre}
\begin{blocLivre}{ex-17-3}\begin{exercice}[1]{Lire un indice}
Un verre a un indice 1,50 et l’air un indice proche de 1,00.
\begin{questions}
\item L’indice a-t-il une unité ?
\item À l’entrée dans le verre, le rayon se rapproche-t-il de la normale ?
\end{questions}
\end{exercice}\end{blocLivre}
\begin{blocLivre}{ex-17-4}\begin{exercice}[1]{Une incidence normale}
Un rayon arrive perpendiculairement à une interface.
\begin{questions}
\item Que vaut $i$ ?
\item Que donne la loi de réfraction pour $r$ ?
\end{questions}
\end{exercice}\end{blocLivre}
\rubrique{S’entraîner}
\begin{blocLivre}{ex-17-5}\begin{exercice}[2]{Calculer un angle réfracté}
De l’air ($n_1=1,00$) vers le verre ($n_2=1,50$), $i=30,0^\circ$.
\schema{refraction}
\begin{questions}
\item Calculer $\sin r$.
\item Déterminer $r$ et contrôler le sens de déviation.
\end{questions}
\end{exercice}\end{blocLivre}
\begin{blocLivre}{ex-17-6}\begin{exercice}[2]{Retrouver un indice}
Un rayon venant de l’air donne $i=45,0^\circ$ et $r=28,1^\circ$.
\begin{questions}
\item Exprimer l’indice du matériau.
\item Calculer sa valeur approchée.
\end{questions}
\end{exercice}\end{blocLivre}
\begin{blocLivre}{ex-17-7}\begin{exercice}[2]{Lire une droite de sinus}
Le graphique porte $\sin i$ verticalement et $\sin r$ horizontalement pour une entrée depuis l’air.
\schema{snell-graphe}
\begin{questions}
\item Déterminer la pente de la droite.
\item En déduire l’indice du milieu.
\end{questions}
\end{exercice}\end{blocLivre}
\begin{blocLivre}{ex-17-8}\begin{exercice}[2]{Pourquoi un prisme disperse}
Le verre d’un prisme possède un indice légèrement différent pour le rouge et le violet.
\schema{prisme}
\begin{questions}
\item Pourquoi les deux radiations ne sortent-elles pas dans la même direction ?
\item Le prisme transforme-t-il une radiation rouge en radiation violette ?
\end{questions}
\end{exercice}\end{blocLivre}
\rubrique{Raisonner et approfondir}
\begin{blocLivre}{ex-17-9}\begin{exercice}[3]{Une sortie vers l’air}
Un rayon passe d’un verre d’indice 1,50 à l’air avec $i=20,0^\circ$.
\begin{questions}
\item Calculer $\sin r$ puis $r$.
\item Le rayon s’éloigne-t-il de la normale ?
\end{questions}
\end{exercice}\end{blocLivre}
\begin{blocLivre}{ex-17-10}\begin{exercice}[3]{Une pente inversée}
Un groupe trace $\sin r$ en fonction de $\sin i$ et trouve une pente 0,667 pour un passage air-verre.
\begin{questions}
\item Pourquoi cette pente n’est-elle pas directement l’indice du verre ?
\item En déduire l’indice.
\end{questions}
\end{exercice}\end{blocLivre}
\begin{blocLivre}{ex-17-11}\begin{exercice}[3]{Un angle impossible}
Pour un passage d’un milieu d’indice 1,50 vers l’air, un calcul conduit à $\sin r=1,15$.
\begin{questions}
\item Pourquoi aucun angle réfracté réel ne correspond-il ?
\item Que doit-on conclure sur le modèle de rayon transmis pour cette incidence ?
\end{questions}
\end{exercice}\end{blocLivre}
\begin{blocLivre}{ex-17-12}\begin{exercice}[3]{Tester une série}
Pour une entrée depuis l’air, les couples $(\sin r;\sin i)$ mesurés sont $(0,10;0,15)$, $(0,20;0,30)$ et $(0,30;0,46)$.
\begin{questions}
\item Comparer les rapports $\sin i/\sin r$.
\item Une petite différence suffit-elle à rejeter la loi ?
\end{questions}
\end{exercice}\end{blocLivre}

